\begin{proposition}\label{prop:model-maximum-principle}
For \(\theta\)-psh model functions \(\varphi,\psi\),
\begin{equation}
\label{eq:model-maximum-principle}
 \one_{\{\varphi>\psi\}}
 \MA_\theta(\max\{\varphi,\psi\})
 =\one_{\{\varphi>\psi\}}\MA_\theta(\varphi).
\end{equation}
\end{proposition}

\begin{proof}
The inputs may be real model functions.  Their maximum
\(\sigma:=\max\{\varphi,\psi\}\) is continuous and \(\theta\)-psh,
but need not be a model function.  Its measure is the bounded product of
\Cref{thm:bounded-ma-extension}, restricted to continuous potentials as
in \Cref{prop:ch24-continuous-uniform}.  We first prove the formula when
the two potentials are rational; the reference form \(\theta\) remains
the given real form throughout.

\emph{Step 1: rational inputs on one common SNC model.}
Suppose \(\varphi,\psi\) are rational model functions.  Then \(\sigma\)
is a rational \(\theta\)-psh model function by
\Cref{prop:ch20-model-max}.  Choose an SNC model \(\cX\) determining
\(\theta\) and all three potentials, using \Cref{prop:snc-cofinal}.
By \Cref{lem:model-function-affine}, all three factor through
\(p_{\cX}\) and are affine on every face of \(\Delta_{\cX}\).
The two measures are supported at its vertices by
\eqref{eq:model-mixed-measure}.  We therefore compare their coefficients
at each \(x_E\) with \(\varphi(x_E)>\psi(x_E)\).

If another component \(F\) meets \(E\), consider an edge
\([x_E,x_F]\).  The restrictions of \(\varphi\), \(\psi\), and their
maximum \(\sigma\) are all affine.  The maximum of two affine functions
on a segment cannot be affine if their strict order reverses along that
segment.  Since \(\varphi(x_E)>\psi(x_E)\), it follows that
\(\varphi(x_F)\geq\psi(x_F)\) and \(\sigma=\varphi\) on the whole
edge.  Write \(\sigma=\varphi_{D_\sigma}\) and
\(\varphi=\varphi_D\), with vertical rational Cartier divisors on
\(\cX\).  The coefficient of any component \(G\) in \(D\) is
\(b_G\varphi(x_G)\), by
\Cref{prop:metrics-divisor-model-function}; the same formula holds for
\(D_\sigma\).  Thus \(D_\sigma-D\) has zero coefficient along \(E\)
and along every component meeting \(E\).  Its support is disjoint from
\(E\).

Choose a positive integer \(m\) clearing the denominators of these
divisors.  The Cartier divisor \(m(D_\sigma-D)\) vanishes on a
neighborhood of \(E\), so its associated line bundle restricts trivially
to \(E\).  Dividing its numerical class by \(m\) gives
\([D_\sigma-D]|_E=0\) in \(N^1(E)_{\mathbf R}\).  Consequently
\[
 (\theta+\ddc\sigma)|_E=(\theta+\ddc\varphi)|_E,
\]
and the two coefficients
\(b_E\bigl((\theta+\ddc\sigma)|_E\bigr)^n\) and
\(b_E\bigl((\theta+\ddc\varphi)|_E\bigr)^n\) in
\eqref{eq:model-mixed-measure} agree.  Summing over all strict vertices
proves \eqref{eq:model-maximum-principle} for rational inputs, without
any rationality assumption on \(\theta\).

\emph{Step 2: approximation at the fixed real form.}
For real \(\theta\)-psh model inputs, apply
\Cref{lem:ch20-fixed-form-rational-approximation} to obtain rational
\(\theta\)-psh model functions \(\varphi_m,\psi_m\) with
\[
 \varepsilon_m:=\max\{\|\varphi_m-\varphi\|_\infty,
                         \|\psi_m-\psi\|_\infty\}\longrightarrow0.
\]
The ample generic class needed for this approximation is part of
\eqref{eq:ch24-standing}; the form itself is not perturbed.
Put \(\sigma_m=\max\{\varphi_m,\psi_m\}\), a rational
\(\theta\)-psh model function.  The elementary maximum inequality gives
\(\|\sigma_m-\sigma\|_\infty\leq\varepsilon_m\).
In particular \(\sigma\) is \(\theta\)-psh by uniform stability,
\Cref{thm:ch21-stability}\,\cref{itm:ch21-stab:unif}.
By \Cref{prop:ch24-continuous-uniform}\,\cref{itm:ch24-cu:unif},
\[
 \MA_\theta(\sigma_m)\longrightarrow\MA_\theta(\sigma),
 \qquad
 \MA_\theta(\varphi_m)\longrightarrow\MA_\theta(\varphi)
 \quad\text{weakly}.
\]

\emph{Step 3: localization before taking the limit.}
The set \(U=\{\varphi>\psi\}\) is open.  Let \(h\) be a continuous
function on \(X^{\An}\) whose support \(K\) is a compact subset of
\(U\).  If \(K\) is nonempty, continuity and compactness give
\(\delta:=\min_K(\varphi-\psi)>0\).  For every sufficiently large
\(m\), we have \(2\varepsilon_m<\delta\), hence
\(K\subseteq\{\varphi_m>\psi_m\}\).  The rational formula from
Step~1 therefore gives
\[
 \int h\,\MA_\theta(\sigma_m)
       =\int h\,\MA_\theta(\varphi_m).
\]
The equality is automatic if \(K\) is empty.  Weak convergence from
Step~2 now yields
\(\int h\,\MA_\theta(\sigma)=\int h\,\MA_\theta(\varphi)\).
The open subspace \(U\) of the compact Hausdorff space \(X^{\An}\)
is locally compact Hausdorff.  Every continuous compactly supported
function on \(U\) extends by zero to a continuous function of the type
just used.  Uniqueness of Radon measures from these tests shows that the
two measures restricted to \(U\) agree, proving
\eqref{eq:model-maximum-principle}.  We have not passed masses of moving
strict-inequality sets through a weak limit.
\end{proof}
